\documentclass[10pt]{beamer}
\usetheme{Madrid}
\usepackage{booktabs}
\usepackage{graphicx}
\usepackage[T1]{fontenc}
\setbeamertemplate{navigation symbols}{}

% ======================= EDIT YOUR DETAILS =======================
\newcommand{\studentname}{Aflah Muhammed P}
% Change the university roll/registration number:
\newcommand{\rollno}{TVE23CSXXX}
% Change your class roll number:
\newcommand{\rollclass}{Roll No: XX}
% Change your guide name:
\newcommand{\guidename}{Prof. Guide Name}
% Change department / college / short-institute if needed:
\newcommand{\deptname}{Department of Computer Science and Engineering}
\newcommand{\collegename}{College of Engineering Trivandrum}
\newcommand{\shortinst}{CET}
% Change the seminar date:
\newcommand{\seminardate}{\today}
% =================================================================

\title[B.Tech Seminar]{ChatInject: How Forged Chat Formatting Tricks AI Agents into Obeying Attackers}
\subtitle{Based on: ChatInject: Abusing Chat Templates for Prompt Injection in LLM Agents (arXiv:2509.22830, 2025)}
\author[\studentname\ (\shortinst)]{\studentname \\ \small \rollno \\ \small \rollclass \\ \small Guide: \guidename}
\institute[]{\deptname \\ \collegename}
\date{\seminardate}

\begin{document}

\begin{frame}[plain]
  \titlepage
\end{frame}

\begin{frame}{Table of Contents}
  \tableofcontents
\end{frame}

\section{Introduction}
\begin{frame}{Introduction}
\begin{itemize}
  \item LLM agents now read email, web pages, and tool outputs to act
  \item This external data creates new prompt-injection attack surfaces
  \item Indirect injection hides commands inside content agents read
  \item Prior attacks use plain text, which models increasingly resist
  \item ChatInject instead forges the model's own chat template
\end{itemize}
\end{frame}

\section{Literature Survey}
\begin{frame}{Literature Survey / Background}
  \small
  \begin{tabular}{@{}p{2.7cm}p{2.3cm}p{5.6cm}@{}}
    \toprule
    \textbf{Title} & \textbf{Author} & \textbf{Summary} \\
    \midrule
    AgentDojo & Debenedetti et al., 2024 & A benchmark environment for evaluating prompt-injection attacks and defenses on tool-using LLM agents in realistic tasks. \\
    \addlinespace
    InjecAgent & Zhan et al., 2024 & A benchmark measuring indirect prompt injection against tool-integrated LLM agents across many attacker scenarios. \\
    \addlinespace
    AutoHijacker & Liu et al., 2025 & An automated indirect prompt-injection attack on LLM agents, cited as related adaptive attack work. \\
    \addlinespace
    \bottomrule
  \end{tabular}
\end{frame}

\section{Motivation}
\begin{frame}{Motivation}
\begin{itemize}
  \item Plain-text injections are increasingly blocked by model training
  \item Chat templates are trusted but rarely studied as an attack surface
  \item Agents struggle to separate read data from real instructions
  \item Multi-turn persuasion is underexplored in injection attacks
\end{itemize}
\end{frame}

\section{Problem Statement}
\begin{frame}{Problem Statement}
  \begin{block}{}
  LLM agents rely on structured chat templates to tell instructions apart from data, but if attacker-controlled external content can forge those templates, the agent may execute injected instructions as if they were legitimate. The paper investigates how severe this template-based vulnerability is and whether current defenses can stop it.
  \end{block}
\end{frame}

\section{Objectives}
\begin{frame}{Objectives}
\begin{itemize}
  \item Expose chat templates as a real prompt-injection attack surface
  \item Design ChatInject to forge native template formatting
  \item Build a Multi-turn persuasion variant for suspicious actions
  \item Measure success across frontier models and two benchmarks
  \item Test whether existing prompt-based defenses hold up
\end{itemize}
\end{frame}

\section{Methodology}
\begin{frame}[allowframebreaks]{Methodology / Approach}
\begin{itemize}
  \item Identify the special role markers used by chat templates
  \item Craft payloads that mimic system, user, and assistant turns
  \item Inject forged turns via external tool output or documents
  \item Add fake assistant replies to build persuasive context
  \item Extend to a Multi-turn dialogue that primes compliance
  \item Evaluate on the AgentDojo and InjecAgent benchmarks
\end{itemize}
\end{frame}

\section{Key Concepts}
\begin{frame}[allowframebreaks]{Key Concepts}
  \begin{description}
    \item[Chat template mimicry] Formatting injected text with real role markers so it looks like a genuine conversation turn.
    \item[Indirect prompt injection] Hiding malicious instructions inside external content the agent reads, such as emails or tool outputs.
    \item[Multi-turn persuasion] Faking a gradual dialogue, including assistant agreement, to normalize and trigger harmful actions.
    \item[Cross-model transferability] One forged-template payload works across many models, even closed-source ones with hidden templates.
    \item[Defense evasion] Prompt-based defenses like spotlighting and sandwiching fail to block the forged-template attack.
  \end{description}
\end{frame}

\section{System Architecture}
\begin{frame}{System Architecture / How It Works}
  \begin{block}{Overview}
  A user gives the agent a task, and the agent calls an external tool or reads a document whose content includes attacker-controlled text. That attacker text is formatted with forged chat-template markers so it appears as a new user or system turn, or as a fake multi-turn exchange with agreeable assistant replies. When this combined text is fed back into the LLM, the model parses the forged turns as legitimate instructions and performs the attacker's action. A diagram can show the loop: User to Agent to Tool or Environment (poisoned) and back into the LLM context, with the forged template block highlighted in red.
  \end{block}
  % TIP: insert a diagram here, e.g.:
  % \begin{center}\includegraphics[width=0.7\textwidth]{your-figure.png}\end{center}
\end{frame}

\section{Results}
\begin{frame}[allowframebreaks]{Results / Findings}
\begin{itemize}
  \item AgentDojo average attack success rate rose from 5.18\% to 32.05\%
  \item InjecAgent average attack success rate rose from 15.13\% to 45.90\%
  \item Multi-turn variant reached 52.33\% success rate on InjecAgent
  \item Forged payloads transferred across models, including closed-source ones
  \item Four prompt-based defenses largely failed, worst against Multi-turn
\end{itemize}
\end{frame}

\section{Discussion}
\begin{frame}{Discussion}
\begin{itemize}
  \item Models cannot reliably separate read data from real instructions
  \item Forged formatting beats plain-text injection by wide margins
  \item Prompt-only defenses can give a false sense of security
  \item Risk grows as agents gain more tool and web access
\end{itemize}
\end{frame}

\section{Challenges}
\begin{frame}{Limitations \& Challenges}
\begin{itemize}
  \item Success rates vary widely by model and benchmark
  \item Evaluation focuses on prompt-based, not structural, defenses
  \item Closed-source template details are inferred, not fully known
  \item Results are tied to two benchmarks and the tested model set
\end{itemize}
\end{frame}

\section{Future Work}
\begin{frame}{Future Work}
\begin{itemize}
  \item Develop structural defenses separating trusted and untrusted input
  \item Train models to resist forged template markers
  \item Extend testing to more agents and real deployments
  \item Study automatic detection of multi-turn persuasion patterns
\end{itemize}
\end{frame}

\section{Conclusion}
\begin{frame}{Conclusion}
\begin{itemize}
  \item Chat templates are a powerful, overlooked attack surface
  \item ChatInject far outperforms plain-text prompt injection
  \item The attack transfers across models and evades prompt defenses
  \item Safer agents need structural, not just prompt-based, protections
\end{itemize}
\end{frame}

\section{References}
\begin{frame}[allowframebreaks]{References}
  \footnotesize
  \begin{enumerate}
    \item ChatInject: Abusing Chat Templates for Prompt Injection in LLM Agents, Hwan Chang, Yonghyun Jun, Hwanhee Lee, arXiv:2509.22830 (ICLR 2026), 2025
    \item AgentDojo: A Dynamic Environment to Evaluate Prompt Injection Attacks and Defenses for LLM Agents, Debenedetti et al., NeurIPS Datasets and Benchmarks, 2024
    \item InjecAgent: Benchmarking Indirect Prompt Injections in Tool-Integrated Large Language Model Agents, Zhan et al., Findings of ACL, 2024
    \item AutoHijacker: Automatic Indirect Prompt Injection Against LLM Agents, Liu et al., 2025
    \item Adaptive Attacks on Prompt Injection Defenses for LLM Agents, Zhan et al., 2025
  \end{enumerate}
\end{frame}

\begin{frame}[plain]
  \begin{center}
    {\Huge Thank You}\\[1em]
    {\large Questions?}
  \end{center}
\end{frame}

\end{document}